% ReLU itself has no trainable parameters; equal axis units show the two slopes.
\par\noindent\begin{minipage}{\linewidth}
\tagstructbegin{tag=Figure,alttext={ReLU 的输入输出曲线，横轴为输入 t，纵轴为输出 ReLU(t)，两轴使用相同的坐标单位。输入小于等于零时输出为零，曲线沿横轴水平延伸；输入大于零时输出等于输入，曲线斜率为一，在原点形成折角。标出的两个例子为负二、零和一、一，即负二变成零，一保留为一。}}
\begin{center}
\begin{tikzpicture}
\begin{axis}[
 width=110mm,height=65mm,axis equal image,
 xmin=-2.6,xmax=2.5,ymin=-.55,ymax=2.6,
 axis lines=middle,xlabel={$t$},ylabel={$\relu(t)$},
 xtick={-2,-1,0,1,2},ytick={1,2},
 tick label style={font=\footnotesize},label style={font=\small},clip=false,
]
\addplot[flowcolor,very thick,domain=-2.4:0,samples=2] {0};
\addplot[flowcolor,very thick,domain=0:2.2,samples=2] {x};
\draw[muted,dashed] (axis cs:1,0)--(axis cs:1,1)--(axis cs:0,1);
\addplot[datacolor,only marks,mark=*,mark size=2pt] coordinates {(-2,0) (1,1)};
\node[font=\footnotesize,above] at (axis cs:-2,.1) {$(-2,0)$};
\node[font=\footnotesize,below right] at (axis cs:1,1) {$(1,1)$};
\node[font=\small] at (axis cs:-1.3,.75) {负数归零};
\node[font=\small,anchor=west] at (axis cs:1.75,1.25) {正数保留};
\end{axis}
\end{tikzpicture}
\captionof{figure}{ReLU 在零处形成折角：负输入输出零，正输入保持原值}
\label{fig:relu-function}
\end{center}
\tagstructend
\end{minipage}\par
